This study represents the first step in a comprehensive credibility assessment for a computational model of a simplified medical device geometry involving conjugate heat transfer. Using \texttt{OpenFOAM}'s \texttt{chtMultiRegionFoam} solver, a rigorous verification effort was conducted, including both code and solution verification, to evaluate numerical accuracy and implementation. 

The results demonstrate that the solver reliably captures laminar flow and heat transfer behavior, with discretization errors from spatial resolution, timestep size, and solver tolerance systematically reduced to acceptable levels. Importantly, this study shows that the choice of resolution directly impacts both numerical accuracy and computational cost, requiring a careful balance between the two. A key finding is that spatial and temporal discretization errors in transient coupled simulations cannot be assessed independently. Co-refinement of mesh and timestep is necessary to obtain reliable estimates of numerical error, while independent refinement can lead to misleading convergence behavior. This highlights the importance of consistent and systematic verification practices when evaluating multiphysics models. The use of a simplified geometry enabled rigorous verification while preserving the dominant fluid and thermal transport physics of the device. This approach demonstrates that meaningful verification can be achieved even when analytical solutions are limited and full system complexity is not yet incorporated. 

The interface energy consistency analysis further supports the coupled solution. The interface imbalance did not exhibit monotonic convergence with mesh refinement, it remained small, approximately $\mathcal{O}(10^{-4})$ of the applied heat source, thus the energy is well conserved across the fluid–solid interface and that the observed convergence behavior is not associated with a non-physical coupling error. This work also emphasizes practical considerations in verification, including computational cost and reproducibility. High-fidelity simulations require significant resources and generate large datasets, particularly for fine mesh and timestep combinations. By identifying an efficient resolution strategy and providing reproducible workflows and publicly accessible artifacts, this study supports broader accessibility and transparency in computational modeling.

This study is aligned with the ASME V\&V 40 framework and contributes a reproducible, risk-informed case study for building model credibility in medical device applications. For the medium–low risk context considered here, the verification approach demonstrates that a level of rigor beyond minimal practice can be achieved through targeted model simplification and systematic analysis. Establishing numerical correctness and quantifying discretization error provides a critical foundation for subsequent validation and uncertainty quantification efforts. Future work will build upon this verification foundation by incorporating validation and uncertainty quantification to complete the full VVUQ framework. Overall, this work demonstrates that systematic verification of coupled fluid–thermal models, even in the absence of analytical solutions for the full system, can provide credible and decision-relevant modeling evidence under practical computational constraints. 